\section{Related Work}
\label{sec:RelatedWork}

\smallTitle{Quantum routing and online path selection}
Quantum networks distribute entanglement through repeaters and end nodes to support long-distance communication, distributed quantum computing, and sensing~\cite{wehner2018quantum,kimble2008quantum}. Entanglement routing differs from classical packet routing because path quality depends on probabilistic entanglement generation and swapping, decoherence, fidelity loss, and limited qubit memory~\cite{briegel1998quantum,dahlberg2021netsquid,bennett1993teleporting,zukowski1993event}. Prior quantum-routing work studies fidelity-aware path selection, multipath/inter-domain routing, benchmarking-based path estimation, resource-constrained routing, and adaptive or learning-based route selection~\cite{wang2025learning,li2025multipath,liu2024qbgp,wang2024adaptive,chaudhary2023quantum,jallowkhan2025adaptive,clayton2024quarc,cicconetti2024scalable,kumar2024routing,leone2021costvector}. These studies establish the importance of online path choice, but they often differ in topology, allocator assumptions, feedback model, and disruption setting, making robustness claims difficult to compare directly.

\smallTitle{Bandit policies for routing decisions}
Bandit methods provide a natural online-learning interface for quantum routing because each decision reveals partial feedback about the selected path, allocation, or context. We compare stochastic, contextual/neural, adversarial, predictive, and hybrid bandit policies rather than treating ``online learning'' as a single category. UCB and Thompson Sampling motivate stochastic baselines~\cite{auer2002finite,thompson1933likelihood}; EXP3 motivates adversarial baselines~\cite{auer2002nonstochastic}; LinUCB, NeuralUCB, and NeuralTS motivate contextual and nonlinear reward models~\cite{li2010contextual,chu2011contextual,zhou2020neuralucb,zhang2022neuralts}; pursuit-style learning automata motivate structured exploitation~\cite{thathachar2011networks}; and informed/contextual bandits motivate forecast-augmented policies~\cite{kar2024icmab,box2015time}. Huang \etal~\cite{huang2024quantum} are closest to our setting through \emph{EXPNeuralUCB}, which combines EXP3-style exploration with NeuralUCB-style reward modeling for joint path selection and qubit allocation.

\smallTitle{Benchmarking and matched robustness evaluation}
Quantum-network benchmarking and utility-driven simulation frameworks motivate separating network behavior from policy behavior during evaluation~\cite{coopmans2021benchmark,kozlowski2022utility}, while best-of-both-worlds bandit theory motivates testing across both stochastic and adversarial regimes~\cite{zimmert2019optimal}. Link-level and path-level best-arm formulations such as LinkSelFiE and related benchmarking approaches address unknown link or path qualities~\cite{10621263,wang2025learning}, but they do not jointly test allocator policy, replay-capacity semantics, and threat progression. Our evaluation targets this matched-comparison gap: all bandit-policy families are tested under the same topology, threat regimes, allocator policies, and replay/capacity settings so that robustness differences can be attributed to the bandit-policy--allocator--capacity interaction rather than to incompatible experimental assumptions.
